\chapter{Analisi dei Requisiti}
\section{Attori e Casi d'Uso}
A scopo di contestualizzazione del progetto è necessario presentare i principali Attori che interagiscono con il Server MCP e i principali scenari d'uso:

\subsection{Attori}
\begin{itemize}
    \item \textbf{Hypervisor bare-metal:} Ambiente di virtualizzazione installato direttamente sull'hardware con cui si ha necessità di interagire.
    
    \item \textbf{Client LLM:} Applicazione client su cui è eseguita l'interfaccia del Large Language Model incaricato di tradurre le richieste dell'utente in linguaggio naturale a payload JSON-RPC.
    
    \item \textbf{Utente (Amministratore di Sistema):} Utente che deve eseguire transazioni di gestione e manutenzione sull'hypervisor tramite il linguaggio naturale.
\end{itemize}

\subsection{Casi d'Uso Principali}
I casi d'uso descritti sono affrontati con granularità SEA-LEVEL.
\begin{itemize}
    \item \textbf{UC1 - Ispezione e monitoraggio dell'infrastruttura:}
    L'Utente richiede una panoramica sullo stato operativo o sulle risorse del cluster di virtualizzazione attraverso il Client LLM.
    \begin{itemize}
        \item \textbf{Attori Coinvolti:} Utente (iniziatore), Client LLM (mediatore), Hypervisor (sistema secondario).
        \item \textbf{Precondizioni:} Il Server MCP è attivo, in ascolto sul canale di trasporto designato e autenticato con successo verso l'API REST dell'Hypervisor tramite token segregato.
        \item \textbf{Flusso Principale degli Eventi:}
        \begin{enumerate}
            \item L'Utente formula una richiesta sullo stato di una o più risorse.
            \item L'LLM interpreta la richiesta, seleziona i tool necessari esposti dall'MCP e genera un corrispondente payload JSON-RPC per il Server MCP.
            \item Il Server MCP intercetta e valida la richiesta in JSON-RPC (\textit{RF2}).
            \begin{itemize}
                \item Se il payload JSON-RPC risulta sintatticamente corretto il sistema continua con la gestione della richiesta.
                \item Se il payload JSON-RPC risulta sintatticamente errato il sistema restituisce un messaggio di errore JSON-RPC all'LLM.
            \end{itemize}
            \item Il Server MCP invoca gli endpoint REST dell'Hypervisor necessari a soddisfare la richiesta.
            \item L'Hypervisor elabora la richiesta ed emette la risposta con codice HTTP 200 contenente le metriche grezze.
            \begin{itemize}
                \item Se l'hypervisor risulta irraggiungibile viene restituito un errore di Connection Timeout.
            \end{itemize}
            \item Il Server MCP estrae e normalizza i dati salienti, restituendoli al Client LLM sotto forma di \textit{Observation} strutturata (\textit{RF3, RF4}).
            \item Il Client LLM sintetizza i dati ricevuti e genera la risposta in linguaggio naturale per l'Utente.
        \end{enumerate}
        \item \textbf{Postcondizioni:} Nessuna modifica di stato viene apportata all'infrastruttura; l'operazione di lettura e le relative chiamate API vengono registrate nel log di sistema (\textit{RF5}).
    \end{itemize}

    \item \textbf{UC2 - Gestione del ciclo di vita e provisioning delle istanze (\textit{VM o Container LXC}):}
    L'Utente desidera la creazione di una istanza di VM o Container LXC specifica.
    \begin{itemize}
        \item \textbf{Attori Coinvolti:} Utente (iniziatore), Client LLM (mediatore), Hypervisor (sistema secondario).
        \item \textbf{Precondizioni:} Il Server MCP è attivo, in ascolto sul canale di trasporto designato e autenticato con successo verso l'API REST dell'Hypervisor tramite token segregato.
        \item \textbf{Flusso Principale degli Eventi:}
        \begin{enumerate}
            \item L'Utente formula la richiesta di voler creare una istanza di VM o Container LXC all'LLM, comunicando nella richiesta anche quale risorse si vogliono allocare.
            \item L'LLM interpreta la richiesta, seleziona i tool necessari esposti dall'MCP e genera un corrispondente payload JSON-RPC per il Server MCP.
            \begin{itemize}
                \item Nel caso non siano state definite tutte le risorse necessarie, l'LLM provvede a richiederle all'Utente prima di compilare un JSON-RPC.
            \end{itemize}
            \item Il Server MCP intercetta e valida la richiesta in JSON-RPC (\textit{RF2}).
            \begin{itemize}
                \item Se il payload JSON-RPC risulta sintatticamente corretto il sistema continua con la gestione della richiesa.
                \item Se il payload JSON-RPC risulta sintatticamente errato il sistema restituisce un messaggio di errore JSON-RPC all'LLM.
            \end{itemize}
            \item Il Server MCP invoca gli endpoint REST dell'Hypervisor necessari a soddisfare la richiesta.
            \item L'Hypervisor elabora la richiesta ed emette la risposta con codice HTTP 200 contenente l'esito della richiesta.
            \begin{itemize}
                \item Se l'hypervisor risulta irraggiungibile viene restituito un errore di Connection Timeout.
            \end{itemize}
            \item Il Server MCP estrae e normalizza l'esito dell'operazione, restituendolo al Client LLM sotto forma di \textit{Observation} strutturata (\textit{RF3, RF4}).
            \item Il Client LLM sintetizza i dati ricevuti e genera la risposta in linguaggio naturale per l'Utente.
        \end{enumerate}
        \item \textbf{Postcondizioni:} Le altre istanze di virtualizzazione non vengono alterate, le relative chiamate API vengono registrate nel log di sistema (\textit{RF5}).
    \end{itemize}
    
    \item \textbf{UC3 - Esecuzione controllata di operazioni critiche (\textit{Human-in-the-Loop}):}
    L'Utente desidera l'eliminazione di una certa istanza di virtualizzazione (Figura \ref{fig:sequence_diagram_uc3}).
    \begin{itemize}
        \item \textbf{Attori Coinvolti:} Utente (iniziatore), Client LLM (mediatore), Hypervisor (sistema secondario).
        \item \textbf{Precondizioni:} Il Server MCP è attivo, in ascolto sul canale di trasporto designato e autenticato con successo verso l'API REST dell'Hypervisor tramite token segregato.
        \item \textbf{Flusso Principale degli Eventi:}
        \begin{enumerate}
            \item L'Utente formula la richiesta di voler eliminare una certa istanza di virtualizzazione.
            \item L'LLM interpreta la richiesta, seleziona i tool necessari esposti dall'MCP e genera un corrispondente payload JSON-RPC per il Server MCP.
            \item Il Server MCP intercetta e valida la richiesta in JSON-RPC (\textit{RF2}), ma riconosce in essa la presenza di comandi critici per soddisfarla (\textit{RF11}).
            \item Il Server sospende l'esecuzione dell'API verso l'Hypervisor e restituisce all'LLM una richiesta formale di validazione, impedendo l'azione silente (\textit{RF6}).
            \item L'LLM informa l'Utente che è necessaria una conferma esplicita e restituisce un riepilogo dell'istanza che verrà eliminata, al fronte di prevedere errori umani.
            \item L'Utente esamina il riepilogo ed eventualmente conferma la richiesta.
            \item L'LLM inoltra la conferma al Server MCP in attesa.
            \item Il Server MCP una volta validata la conferma procede con l'inoltro dei comandi verso l'Hypervisor.
            \item L'Hypervisor effettua l'eliminazione e restituisce l'esito al Server MCP.
            \item Il Server MCP comunica l'effettiva eliminazione all'LLM che a sua volta notifica l'utente.
        \end{enumerate}
        \item \textbf{Postcondizioni:} L'istanza interessata viene eliminata, le istanze di virtualizzazione non interessate non vengono alterate, le relative chiamate API inclusa la conferma vengono registrate nel log di sistema (\textit{RF5}).
    \end{itemize}
\end{itemize}

\section{Requisiti Funzionali}
\subsection{Interfacciamento con Hypervisor e Protocollo MCP}
\begin{itemize}
    \item \textbf{RF1 - Esposizione dinamica dei tool (\textit{Tool Discovery}):} Il layer MCP deve esporre in modo dinamico i tool per le interazioni con l'infrastruttura IT in formato \textit{JSON Schema}.
    
    \item \textbf{RF2 - Layer di traduzione (JSON-RPC verso REST):} Il layer MCP si occupa di tradurre dei payload JSON-RPC generati dall'LLM e controllare se conformi ai suoi standard per poi inoltrare le eventuali richieste REST API all'hypervisor.
    
    \item \textbf{RF3 - Formattazione dei risultati per il contesto del LLM client:} Il server MCP dovrà restituire un output formattato in modo strutturato e fruibile per il client LLM.
\end{itemize}

\subsection{Ispezione e Monitoraggio dell'Infrastruttura (Read-Only)}
\begin{itemize}
    \item \textbf{RF4 - Monitoraggio dello stato hardware e logico del sistema:} Il sistema deve offrire la possibilità di ricevere una panoramica completa del sistema in termini di risorse allocate ed occupate in tempo reale dai nodi nonchè dalle VM e Container in servizio.
    
    \item \textbf{RF5 - Tracciamento delle transazioni (\textit{Audit Logging}):} Il layer MCP deve mantenere uno storico di tutte le richieste ricevute in termini di utente, data e ora, JSON-RPC, API invocate per soddisfare la richiesta e codice HTTP ricevuto.
    
    \item \textbf{RF6 - Sistema di Notifica:} All'atto di una richiesta che necessita intervento umano, il sistema deve inviare un alert all'utente interessato.
\end{itemize}

\subsection{Gestione del Ciclo di Vita delle Istanze (Operazioni di Stato)}
\begin{itemize}
    \item \textbf{RF7 - Controllo dell'alimentazione (\textit{Power Management}):} Devono essere esposti strumenti di controllo dell'alimentazione di VM e Container in modo sicuro (\textit{boot, shutdown, reboot}).
    
    \item \textbf{RF8 - Provisioning di base:} Devono essere esposti strumenti per la creazione, gestione dinamica delle risorse allocate e distruzione controllata di VM e Container, anche tramite template registrati sull'hypervisor.
    
    \item \textbf{RF9 - Gestione degli snapshot:} Prima di ogni intervento di manutenzione deve essere consigliata dall'LLM la possibilità di eseguire uno snapshot del sistema, consentendo un eventuale ripristino \textit{rollback} all'occorrenza.
\end{itemize}

\subsection{Sicurezza e Controllo d'Esecuzione}
\begin{itemize}
    \item \textbf{RF10 - Autenticazione presso l'hypervisor e gestione dei permessi:} Il layer MCP deve caricare le credenziali d'accesso dell'utente attraverso variabili d'ambiente per poi mantenere un token di autenticazione valido per interagire con l'hypervisor. Le credenziali non devono transitare tramite il contesto conversazionale dell'LLM.
    
    \item \textbf{RF11 - Intercettazione e blocco di comandi distruttivi:} Il sistema deve saper intercettare e richiedere conferma all'utente nel caso certi prompt prevedano comandi possibilmente distruttivi o che possano apportare modifiche architetturali importanti.
\end{itemize}

\section{Requisiti Non Funzionali}
\begin{itemize}
    \item \textbf{RNF1 - Modularità ed Estendibilità:} Deve essere possibile interfacciare diversi LLM con diversi Hypervisor e prevedere un'architettura tale da rendere semplice aggiungere o deprecare funzionalità dell'MCP Server.

    \item \textbf{RNF2 - Efficienza delle prestazioni (latenza):} La piattaforma deve introdurre il minimo \textit{overhead} computazionale e garantire che un certo prompt esegua i comandi sul sistema e riceva un output in un tempo trascurabile.
    
    \item \textbf{RNF3 - Disponibilità (\textit{Availability}):} La piattaforma deve garantire un alto livello di disponibilità anche sotto elevato stress di carico.
    
    \item \textbf{RNF4 - Concorrenza:} La piattaforma deve prevedere paradigmi di gestione delle concorrenze e operzioni asincrone non bloccanti.

    \item \textbf{RNF5 - Tolleranza ai guasti (\textit{Fault Tolerance}:} L'architettura deve essere tale da implementare una gestione robusta degli imprevisti, quali possibili timeout di risposta, errori o JSON malformati e incoerenti.
    
    \item \textbf{RNF6 - Consistenza e Pattern Saga:} Logicamente ove un prompt preveda più comandi, esso deve essere considerato eseguito solo se tutti i comandi hanno successo, ma in termini pratici diventa ingegneristicamente complesso trasformare un prompt in una trasizione (\textit{ACID}). Per questo motivo, bisogna quindi che il sistema interrompa eventuali sequenze di comandi ne caso uno dovesse fallire e proporre eventualmente delle \textit{operazioni di rollback}.

    \item \textbf{RNF7 - Sicurezza e supervisione (\textit{Human-in-the-Loop}):} Bisogna che azioni critiche vengano monitorate dall'utente in tempo reale e all'occorenza validate da esso. Tutte le azioni possibilmente distruttive o che possano apportare modifiche profonde al sistema devono ricevere conferma da un utente umano. Deve inoltre essere previsto, quando possibile, il tempo di down del sistema su alcuni comandi complessi ed esporre la scelta all'utente.

    \item \textbf{RNF8 - Manutenibilità:} Il codice deve essere strutturato affinchè in caso di manutenzione il sistema rimanga in down in tempi brevi. Deve essere implementato un sistema di logging granulare sullo storico di ogni prompt per utente e sulle chiamate API che ne sono conseguite in modo da avere il controllo completo e documentato su tutte le azioni dell'LLM.
\end{itemize}